\subsection{A forty-arrow category of type $\mathbb CP^2$}
\label{C:CP2-upper40-section}

\begin{proposition}\label{C:CP2-upper40}
There is a category with two objects and forty arrows whose
classifying space is homotopy equivalent to $\mathbb CP^2$.
Its complete composition table is supplied by
\path{code/d_cp2_category40.py} and
\path{code/outputs/d_cp2_category40_verified.json}.
\end{proposition}
\begin{proof}
We first specify the forty-arrow table directly. Let
$L_0=\{1,x_{00},x_{01},x_{10},x_{11}\}$, with
$x_{ij}x_{kl}=x_{il}$. Put
\[
\begin{split}
 O={}&\{(0,1,2),(0,0,0),(0,0,1),(0,0,2),
             (1,1,1),(1,1,2),(2,2,2)\},\\
 D={}&\{(0,0,0),(0,0,1),(1,1,1),(2,2,2)\}.
\end{split}
\]
These tuples are monotone self-maps of $\{0,1,2\}$.
Let $A$ be the 29-element monoid
\[
 A=(\{1,x_{00},x_{01}\}\times O)
       \cup(\{x_{10},x_{11}\}\times D),
 \qquad
 (l,\theta)(m,\eta)=(lm,\eta\circ\theta).
\]
The identity is $(1,(0,1,2))$. Define the ten-element set
\[
 S=\{(a,b,c):a\in\{0,1\},\ a\leq b\leq c\leq a+2\}
          \setminus\{(1,2,2),(1,3,3)\}.
\]
For an integer tuple $s$ write
$\nu(s)=s-2\lfloor s_0/2\rfloor(1,1,1)$, and define
$\lambda_1(t)=t$ and
$\lambda_{x_{ij}}(t)=2\lfloor(t+j)/2\rfloor+i$.
The left $A$-action on $S$ is
\[
 (l,\theta)\cdot s
   =\nu\bigl(\lambda_l(s_{\theta(0)}),
              \lambda_l(s_{\theta(1)}),
              \lambda_l(s_{\theta(2)})\bigr).
\]
The displayed sets are closed under these operations. The resulting
category $C$ has objects $a,*$, with
\[
 C(a,a)=A,\quad C(*,*)=\{1_*\},\quad
 C(*,a)=S,\quad C(a,*)=\varnothing.
\]
Composition with an arrow in $S$ is exactly the displayed action;
all other compositions are multiplication in $A$ or identities.
Thus the arrow count is $29+1+10=40$.
The independent table check verifies all 1142 typed pairs and
33\,100 associativity triples.

To prove the homotopy type, adjoin an intermediate object $b$ with
$\operatorname{End}(b)=\{1_b\}$, arrows
$i_k:b\to a$ for $k=0,1$, and arrows
$r_{j,t}:a\to b$ for $j=0,1$ and $t=0,1,2$. Specify
\[
 r_{j,t}i_k=1_b,\qquad
 i_kr_{j,t}=(x_{kj},\operatorname{const}_t).
\]
The old elements act by
\[
 (l,\theta)i_k=
 \begin{cases}i_k&l=1,\\i_i&l=x_{ij},\end{cases}
 \qquad
 r_{j,t}(l,\theta)=
 \begin{cases}
 r_{j,\theta(t)}&l=1,\\
 r_{v,\theta(t)}&l=x_{uv}.
 \end{cases}
\]
Extend $S$ to the covariant set-valued functor $P$ on the
two-object subcategory $K=\{a,b\}$ by
$P(a)=S$, $P(b)=\{s_0\}$,
$P(i_k)(s_0)=(k,k,k)$, and $P(r_{j,t})$ the unique map.
The collage $G^-$ of $P$ has fifty arrows: it adds
$*\to b$ to the original category and the nine other arrows
involving $b$. Its full table was independently checked on all
43\,302 composable triples.

The inclusion $C\hookrightarrow G^-$ is a homotopy equivalence.
For a target in $C$ its comma category has a terminal object.
For the target $b$, choose $i_0$ and $r_{0,0}$ with
$r_{0,0}i_0=1_b$. The constant functor to
$(a,r_{0,0})\in(C\downarrow b)$ receives a natural transformation
from the identity: at $(d,\alpha:d\to b)$ its component is
$i_0\alpha:d\to a$. This contracts the comma category.
Quillen's Theorem A proves the assertion.

Since $P(b)$ is a singleton and $b$ retracts from $a$,
Lemma~\ref{lem:B-framed-collage} identifies $BG^-$ with the
classifying space of the explicit monoid
\[
 T=A\amalg\bigl(S\times\{0,1\}\times\{0,1,2\}\bigr),
 \qquad |T|=29+60=89.
\]
The second summand is a rectangular band; the left action is the
one just displayed and the right action is
\[
 (s,j,t)(l,\theta)=
 \begin{cases}
 (s,j,\theta(t))&l=1,\\
 (s,v,\theta(t))&l=x_{uv}.
 \end{cases}
\]
Identify the twelve pairs
\[
 (x_{kj},\operatorname{const}_t)=((k,k,k),j,t).
\]
The resulting quotient is exactly the 77-element monoid of
Proposition~\ref{prop:B-CP77}; the program checks every equation
of this specified quotient map $T\to M_{77}$.

Here the homology-equivalence assertion is established by a complete
resolution and its actual comparison map. Number $T$ by listing
the old elements in the displayed $L_0$ order, each with the
displayed $O$ or $D$ order, followed by the band in lexicographic
row, $j$, $t$ order. Write $u_i$ for element $i$, starting at zero,
and $R=\Z T$. The finite projective resolution has terms
\[
\begin{array}{c|l|r}
 n&P_n&\operatorname{rank}_{\Z}P_n\\\hline
 0&Ru_{29}&10\\
 1&\bigl(R(u_0-u_3)\bigr)^3&108\\
 2&R(u_8-u_{29})\oplus Ru_{29}\oplus
        \bigl(R(u_7-u_{10})\bigr)^5\oplus
        \bigl(R(u_3-u_1)\bigr)^4&204\\
 3&\bigl(R(u_1-u_8)\bigr)^2\oplus
        \bigl(R(u_{10}-u_8)\bigr)^7&117\\
 4&R(u_8-u_{29})\oplus Ru_{29}&12.
\end{array}
\]
Each displayed difference is an idempotent formed from two nested
idempotents. The complete integral differential matrices are in
\path{code/outputs/b_resolution_D_CP50_framed89.json}.
The standalone verifier reconstructs their $T$-actions, verifies
equivariance for all 89 elements, computes saturated images by
integer Smith form, checks every kernel rank and zero composite,
and verifies that the final kernel is zero. After tensoring with
the trivial right module, only one copy of $\Z$ in degrees
zero, two and four remains.

The certificate
\path{code/outputs/b_resolution_map_D_CP50_framed89_to_CP77.json}
is an augmentation-preserving comparison for the specified quotient,
constructed by integral lifts on each idempotent generator. The
standalone comparison verifier checks every chain square and every
source action and recomputes the maps on trivial-coefficient chains.
They are $(1)$ in each of degrees zero, two and four. Hence
$BT\to BM_{77}$ is an integral homology equivalence. Every element
of $T$ has eventually constant powers, as also checked directly,
so its universal group is trivial. The target is likewise simply
connected. The homological Whitehead theorem now gives
\[
 BC\simeq BG^-\simeq BT\simeq BM_{77}\simeq\mathbb CP^2.
\]
This uses the literal comparison map in degree four; equal additive
homology alone, or an isomorphism only in degree two, would not
identify the homotopy type.
\end{proof}

\begin{proposition}[An all-order obstruction to literal small-target Hopf maps]
\label{C:aperiodic-Hopf-target}
\label{C:literal-Hopf-target}
Let $M$ be a finite monoid with $H_2(BM;\Q)=0$.
Every unital homomorphism $M\to L_0$ induces a nullhomotopic map
of classifying spaces. In particular no such homomorphism from
a finite monoid of type $S^3$ represents a Hopf map to the
five-element model of $S^2$.
\end{proposition}
\begin{proof}
First collapse each minimal $\mathcal H$-class of $M$, obtaining
the actual quotient $M\to\overline M$ of
Lemma~\ref{lem:B-collapse-minimal-group}. This is a rational
homology equivalence. Each minimal $\mathcal H$-class is a finite
group with its own idempotent identity, and a homomorphism into a
band is constant on that group. Therefore every map to $L_0$
factors through $\overline M$. We may replace $M$ by this quotient
and suppose its minimal ideal is a rectangular band.

Suppose that $\phi:M\twoheadrightarrow L_0$ is onto.
Its image sends the minimal ideal of $M$ onto the minimal ideal
of $L_0$. Write its rectangular minimal ideal as $K=I\times J$;
the target ideal is the
$2\times2$ band. Choose $e\in K$ and write $f=\phi(e)$.
The maps of principal ideals
\[
 Me\longrightarrow L_0f,\qquad eM\longrightarrow fL_0
\]
are surjections of the corresponding row and column sets.
They induce surjections of their rational augmentation kernels
$V_M\twoheadrightarrow V_{L_0}$ and
$W_M\twoheadrightarrow W_{L_0}$: lift each difference from a
chosen base row or column to a difference in the source.
There is also a surjective ring map
\[
 A_M=\Q M/\Q K\longrightarrow
 A_{L_0}=\Q L_0/\Q(K_{L_0}).
\]

The two connecting maps used in Lemma~\ref{lem:B-band-Tor}
are natural for this homomorphism. Indeed, apply $\phi$ to
the two augmentation sequences for $Me,eM$, and to their
two-sided bar complexes; this gives commuting maps on both
connecting homomorphisms. Under the resulting natural
identification of degree-two homology, the induced map is
therefore exactly
\[
 W_M\otimes_{A_M}V_M\longrightarrow
 W_{L_0}\otimes_{A_{L_0}}V_{L_0},
 \qquad w\otimes v\longmapsto\phi(w)\otimes\phi(v).
\]
It is onto, since every pure tensor has lifts in both factors.
Thus $H_2(BM;\Q)\to H_2(BL_0;\Q)=\Q$ is onto, contradicting
the hypothesis.

Consequently the image of $\phi$ is a proper unital submonoid
of $L_0$. If its band elements use two different rows and two
different columns, products fill the entire $2\times2$ band,
so the image would be all of $L_0$. Its nonidentity elements
therefore lie in a single row or a single column. If this set
is nonempty, any of its elements is respectively a right zero
or a left zero in the image monoid. Such a zero gives a natural
transformation between the identity functor and the constant
functor, so the image has contractible classifying space.
The image $\{1\}$ is also contractible. The map $B\phi$
factors through that contractible space, proving the claim.
\end{proof}

\begin{remark}[Scope of subsequent reductions]\label{C:CP40-reductions}
The program \texttt{c\_CP40\_all\_congruences.py} enumerates the
complete lattice of 28\,135 congruences of the forty-arrow table
that preserve its two objects and hom-set types. Among the proper
quotients, 25\,704 have zero rational degree-two homology at the
endomorphism-monoid vertex, which forces the degree-two square
of the act cone to vanish; the other 2430 have zero degree-two
homology over $\mathbb F_2$. Thus no proper quotient in this
specified sense is another $\mathbb CP^2$ model.

There are ten maximal proper wide subcategories of this table.
Their complete act-cone projective resolutions, reconstructed and
checked independently, all have different additive homology from
$\mathbb CP^2$. This does not exclude more deeply nested
subcategories. A separate search of all sixteen smaller full
subcategories retaining the cone object and at least one of its
proper Karoubi-object classes finds either the wrong second Betti
number or a zero degree-two square. These are explicitly scoped
reduction checks, not a proof that forty is the category minimum.
The records are
\path{code/outputs/c_CP40_all_congruences.json},
\path{code/outputs/d_cp2_subcategory_search.json}, and
\path{code/outputs/d_cp2_karoubi_search.json}.
\end{remark}
